\documentclass[10pt,a4paper]{amsart}
\usepackage[margin=19mm]{geometry}
\usepackage{amsmath,amssymb,mathtools}
\usepackage[scr=rsfs,cal=euler]{mathalfa}
\usepackage{xfrac}
\usepackage[unicode,hidelinks,bookmarks=false]{hyperref}
\newcommand{\ie}{i.e.\ }
\newcommand{\cf}{cf.\ }
\newcommand{\resp}{respectively\ }
\newcommand{\Subsection}{\subsection}
\providecommand{\nobreakdash}{\nobreak\hbox{-}}
\setlength{\emergencystretch}{2em}
\multlinegap0pt

\theoremstyle{plain}
\newtheorem*{mthm}{Main Theorem}
\newtheorem{thm}{Theorem}
\newtheorem{lem}{Lemma}
\newtheorem{cor}{Corollary}
\newtheorem{prop}{Proposition}
\newtheorem{rmk}{Remark}
\newtheorem{defn}{Definition}
\title[On the monomials of Poincar\'e Series of negative index]{On the monomials of Poincar\'e Series of negative index}
\author{Divyanshu Kala}
\date{Source: arXiv:2609.28350v1 (CC BY 4.0)}
\begin{document}
\maketitle

\noindent\emph{Source and licence.} The delivered work is the article shown in the title above, version arXiv:2609.28350v1 (CC BY 4.0), distributed under the Creative Commons Attribution 4.0 International licence, as recorded in the arXiv licence metadata for this exact version and shown on the abstract page saved with this item. The full licence notice, which carries the licence URL and the address of that abstract page, is the bundled file \texttt{licenses/arxiv-2609.28350-CC-BY-4.0.txt}.
\par\smallskip

\noindent\textit{Editorial scope.} Counted unit: the article's Lemma (source label lk>l/2) (source0.tex lines 877--882, source label \texttt{lk>l/2}): ``Let be an even integer such that . For an integer such that , we have l,m (1) - (l-2 l,m ) 2,m (1) 0.02 (l-4m )(l-4m 8) .''. It is delivered together with the article's complete original proof (source0.tex lines 883--924, one proof environment adjacent to the statement, 42 lines), copied line for line from the article's own text. Required context retained uncounted: 0.09 (lines 437--439); k>l/2 (lines 474--477). Editorial formatting only: an AMS \texttt{amsart} preamble that restates the article's own macro definitions and its own theorem declarations, and the theorem environments the retained lines use; the article's own class file is neither shipped nor input; the retained environments are renumbered sequentially in order of appearance, whereas the article prints them with its own numbering; the article's equation numbering is not reproduced and every cross-reference inside the retained text resolves to the wrapper's numbering. No citation occurs inside any retained line, so the wrapper carries no bibliography; All retained bodies are exact copies of the bundled \texttt{sources/211580/source0.tex}, so no omission receipt applies. No asset-inclusion command occurs inside any retained span and the package ships no image asset.


\section*{Retained context: 0.09 (source0.tex lines 437--439)}

\begin{equation}\label{0.09}
0<-m\pi\left(\frac{2\pi}{l-3.96m\pi}\right)^2<0.05.
\end{equation}


\section*{Retained context: k>l/2 (source0.tex lines 474--477)}

\begin{align}\label{k>l/2}
b_{(l-2\lambda_{l,m})/2}(\alpha_{l,m}^{(1)})-b_{(l-2\lambda_{l,m})/2}(\alpha_{(l-2\lambda_{l,m})/2,m})
>\frac{\pi}{2}-\frac{l+2\lambda_{l,m}}{4}\left(\frac{2\pi}{l-4m\pi\left(1-\frac{1}{2}\left(\frac{2\pi}{l-3.96m\pi}\right)^2\right)}\right).
\end{align}


\section{The counted unit: Lemma (source label lk>l/2) and its complete original proof}

\begin{lem}\label{lk>l/2}
Let $l\ge50$ be an even integer such that $r_l=2$. For an integer $m\le-1$ such that $l/2-\lambda_{l,m}\ge 50$, we have
$$
\alpha_{l,m}^{(1)}-\alpha_{(l-2\lambda_{l,m})/2,m}^{(1)}>\frac{0.02\pi}{(l-4m\pi)(l-4m\pi+8)}.
$$
\end{lem}

\begin{proof}
Using \eqref{k>l/2}, we have
\begin{align*}
b_{(l-2\lambda_{l,m})/2}(\alpha_{l,m}^{(1)})-b_{(l-2\lambda_{l,m})/2}(\alpha_{(l-2\lambda_{l,m})/2,m}^{(1)})&>\frac{\pi}{2}-
\frac{l+2\lambda_{l,m}}{4}\left(\frac{2\pi}{l-4m\pi\left(1-\frac{1}{2}\left(\frac{2\pi}{l-2\sqrt3m\pi}\right)^2\right)}\right)\\
&=\frac{\pi}{2}\left(1-\frac{l+2\lambda_{l,m}}{l-4m\pi\left(1-\frac{1}{2}\left(\frac{2\pi}{l-2\sqrt3m\pi}\right)^2\right)}\right).
\end{align*}
As $\lambda_{l,m}\le \lambda_{l,m}'\in J_{l,m}$, we have
\begin{equation}\label{0.005}
\lambda_{l,m}\le -2m\pi\left(1-\frac{1}{2}\left(\frac{2\pi}{l-2\sqrt3m\pi}\right)^2\right).
\end{equation}
For $l\ge 50$ and $m\le-1$, using \eqref{0.09} we have
$$
0<-m\pi\left(\frac{2\pi}{l-2\sqrt3m\pi}\right)^2<0.05
$$
Since $ \{-2m\pi\}\ge 0.06$, we have
$$
\lambda_{l,m}\le -2m\pi\left(1-\frac{1}{2}\left(\frac{2\pi}{l-2\sqrt3m\pi}\right)^2\right)-0.01
$$
Using this, we get that
\begin{align*}
b_{(l-2\lambda_{l,m})/2}(\alpha_{l,m}^{(1)})-b_{(l-2\lambda_{l,m})/2}(\alpha_{(l-2\lambda_{l,m})/2,m}^{(1)})&>\frac{\pi}{2}\left(1-\frac{l-4m\pi\left(1-\frac{1}{2}\left(\frac{2\pi}{l-2\sqrt3m\pi}\right)^2\right)-0.02}{l-4m\pi\left(1-\frac{1}{2}\left(\frac{2\pi}{l-2\sqrt3m\pi}\right)^2\right)}\right)\\
&=\frac{0.01\pi}{l-4m\pi\left(1-\frac{1}{2}\left(\frac{2\pi}{l-2\sqrt3m\pi}\right)^2\right)}.
\end{align*}
Using the mean value theorem, we get that for some $\xi\in (\pi/2,2\pi/3)$,
$$
\alpha_{l,m}^{(1)}-\alpha_{(l-2\lambda_{l,m})/2,m}^{(1)}>\frac{0.01\pi}{\left(l-4m\pi\left(1-\frac{1}{2}\left(\frac{2\pi}{l-2\sqrt3m\pi}\right)^2\right)\right)b_{(l-2\lambda_{l,m})/2}'(\xi)}
$$
Now 
$$
b_{(l-2\lambda_{l,m})/2}'(\xi)\le \frac{l-2\lambda_{l,m}-8m\pi}{4}.
$$
Using this, we get that
$$
\alpha_{l,m}^{(1)}-\alpha_{(l-2\lambda_{l,m})/2,m}^{(1)}>\frac{0.04\pi}{\left(l-4m\pi\left(1-\frac{1}{2}\left(\frac{2\pi}{l-2\sqrt3m\pi}\right)^2\right)\right)(l-2\lambda_{l,m}-8m\pi)}.
$$
As $\lambda_{l,m}>-2m\pi-4$, we get 
\begin{align*}
\alpha_{l,m}^{(1)}-\alpha_{(l-2\lambda_{l,m})/2,m}^{(1)}&>\frac{0.04\pi}{\left(l-4m\pi\left(1-\frac{1}{2}\left(\frac{2\pi}{l-2\sqrt3m\pi}\right)^2\right)\right)(l-4m\pi+8)}\\
&>\frac{0.04\pi}{(l-4m\pi)(l-4m\pi+8)}.
\end{align*}
\end{proof}

\end{document}
